% ===============================================
% ECEV 33400: Stochastic Processes      Fall 2017
% hw_draft.tex
% Template for first draft of homework submission
% ===============================================


% -----------------------------------------------
% The preamble that follows can be ignored. Go on
% down to the section that says "START HERE" 
% -----------------------------------------------

\documentclass{article}

\usepackage[margin=1in]{geometry} 
\usepackage{amsmath,amsthm,amssymb}

\newcommand{\R}{\mathbb{R}}
\newcommand{\Z}{\mathbb{Z}}
\newcommand{\N}{\mathbb{N}}
\newcommand{\Q}{\mathbb{Q}}
\newcommand{\C}{\mathbb{C}}
\newcommand{\F}{{\mathbb F}}
\newcommand{\Pb}{{\mathbb P}}
\newcommand{\E}{{\mathbb E}}
\newcommand{\Frac}[2]{\left(\frac{#1}{#2}\right)}
\newcommand{\Eval}[2]{\big|_{#1}^{#2}}
\newcommand{\intinf}{\int_{-\infty}^{\infty}}

\newenvironment{theorem}[2][Theorem]{\begin{trivlist}
\item[\hskip \labelsep {\bfseries #1}\hskip \labelsep {\bfseries #2.}]}{\end{trivlist}}
\newenvironment{lemma}[2][Lemma]{\begin{trivlist}
\item[\hskip \labelsep {\bfseries #1}\hskip \labelsep {\bfseries #2.}]}{\end{trivlist}}
\newenvironment{exercise}[2][Exercise]{\begin{trivlist}
\item[\hskip \labelsep {\bfseries #1}\hskip \labelsep {\bfseries #2.}]}{\end{trivlist}}
\newenvironment{problem}[2][Problem]{\begin{trivlist}
\item[\hskip \labelsep {\bfseries #1}\hskip \labelsep {\bfseries #2.}]}{\end{trivlist}}
\newenvironment{question}[2][Question]{\begin{trivlist}
\item[\hskip \labelsep {\bfseries #1}\hskip \labelsep {\bfseries #2.}]}{\end{trivlist}}
\newenvironment{corollary}[2][Corollary]{\begin{trivlist}
\item[\hskip \labelsep {\bfseries #1}\hskip \labelsep {\bfseries #2.}]}{\end{trivlist}}

\newenvironment{solution}{\begin{proof}[Solution]}{\end{proof}}


\begin{document}

% ------------------------------------------ %
%                 START HERE                 %
% ------------------------------------------ %

\title{Homework 1} % Replace X with the appropriate number
\author{ \iffalse STUDENT NAME \fi \\ECEV 33400} % Replace  "\iffalse X \fi" with your name.
\date{Due Date: 10/5/2017} % Don't change this! :( 

\maketitle

% -----------------------------------------------------
% The following two environments (theorem, proof) are
% where you will enter the statement and proof of your
% first problem for this assignment.
%
% In the theorem environment, you can replace the word
% "theorem" in the \begin and \end commands with
% "exercise", "problem", "lemma", etc., depending on
% what you are submitting. Replace the "x.yz" with the
% appropriate number for your problem.
%
% If your problem does not involve a formal proof, you
% can change the word "proof" in the \begin and \end
% commands with "solution".
% -----------------------------------------------------

\subsection*{Problems} 
% You can copy these problems down to the environments below to write your solutions & delete the top subsections.

\begin{enumerate}

\item Show that the probability-generating function for the binomial distribution $p_j = \binom{n}{j} p^j q^{n-j}$, where $ q = 1 - p$, is $P(x) = (q + px)^n$. Hence prove that the mean is $m = np$, and the variance is $\sigma^2 = npq$. 

\item What is the probability-generating function for the Poisson distribution $p_j = \frac{e^{-\lambda} \lambda^j}{j!}, j = 0, 1, 2, \cdots$? Hence show that the mean and variance of the distribution are both equal to $\lambda$. 

\item Find the probability-generating function of the geometric distribution $p_j = pq^j, j = 0, 1, 2, \cdots$. What are the mean and variance? 

\item Suppose that $S_N = X_1 + X_2 + \cdots + X_N$, where the $X_j$ are independent random variables with identical distributions given by $$P \{ X_j = 0 \} = q, P \{ X_j = 1 \} = p,$$ and $N$ is a random variable with a Poisson distribution having mean $\lambda$. Prove that $S_N$ has a Poisson distribution with mean $\lambda p$. 
\end{enumerate}

% -----------------------------------------------------
% First problem
% -----------------------------------------------------

% \begin{problem}{1}
% Show that the probability-generating function for the binomial distribution $p_j = \binom{n}{j} p^j q^{n-j}$, where $ q = 1 - p$, is $P(x) = (q + px)^n$. Hence prove that the mean is $m = np$, and the variance is $\sigma^2 = npq$.
% \end{problem}

% \begin{solution}
% \end{solution}

% -----------------------------------------------------
% Second problem
% -----------------------------------------------------

% \vspace{0.25in} % This just adds some space between problems 1 and 2.

% \begin{problem}{2}
% What is the probability-generating function for the Poisson distribution $$p_j = \frac{e^{-\lambda} \lambda^j}{j!}, j = 0, 1, 2, \cdots ?$$ Hence show that the mean and variance of the distribution are both equal to $\lambda$. 
% \end{problem}

% \begin{solution}
% \end{solution}

% -----------------------------------------------------
% Third problem
% -----------------------------------------------------


% -----------------------------------------------------
% Fourth problem
% -----------------------------------------------------


% -----------------------------------------------
% Ignore everything that appears below this.
% -----------------------------------------------

\end{document}
